从早期版本开始，标准库就提供了操作文件的选项。通过流（例如 \texttt{fstream}）可以打开文件、读取数据、写入字节，并执行许多其他操作。然而此前一直缺失的是把文件系统作为一个整体来操作的能力。例如在 C++14 中，你必须借助一些第三方库才能遍历目录、拼接路径、删除目录或读取文件权限。而现在，随着 C++17 的到来，我们以 \texttt{std::filesystem} 组件的形式向前迈进了一大步！

在本章中你将学到：

\begin{itemize}
\item \texttt{std::filesystem} 是如何进入标准的
\item 基本的类型与操作有哪些
\item 如何处理路径
\item 如何处理 \texttt{std::filesystem} 中的错误
\item 如何遍历一个目录
\item 如何创建新的目录与文件
\end{itemize}

\section{文件系统概览}

虽然标准库曾缺少一些重要特性，但你总是可以使用 Boost 及其数十个子库来完成工作。C++ 委员会认为 Boost 库极为重要，于是将其中一些系统合并进了标准。例如智能指针（尽管在 C++11 中借助移动语义得到了改进）、正则表达式、\texttt{std::optional}、\texttt{std::any} 等等。

\texttt{std::filesystem} 也有着相似的故事。

文件系统库直接以 Boost filesystem 为模型，后者自 2003 年（1.30 版本）起就已可用。在 C++ 实现中，委员会还把该组件扩展到了非 POSIX 系统。该库最早以 TS（技术规范，Technical Specification）的形式提供，在经过了长时间的改进与反馈之后，最终合并进了 C++17 标准。

该库位于 \texttt{<filesystem>} 头文件，并使用 \texttt{std::filesystem} 命名空间。

\subsection{库的核心组成}

文件系统库是标准库中相当重要的一部分。它定义了许多带有数十个方法的类型，同时还提供了大量自由函数。

下面我们可以定义该模块的核心元素：

\begin{itemize}
\item \texttt{std::filesystem::path} 对象允许你操纵那些代表系统中存在或不存在的文件与目录的路径。
\item \texttt{std::filesystem::directory\_entry} 代表一个现有的路径，并带有额外的状态信息，例如最后写入时间、文件大小或其他属性。
\item 目录迭代器允许你遍历给定的目录。该库提供了递归与非递归两个版本。
\item 许多辅助函数，例如获取路径信息、文件操作、权限、创建目录等等。
\end{itemize}

在下一节中，你将看到组成 \texttt{std::filesystem} 的所有部分的演示。

\section{演示}

与其一开始就逐个细节地探究该库，不如在下一页先看一个演示示例：递归地显示给定目录中所有文件的基本信息。这能让你对该库的整体样貌有一个高层次的了解。

\begin{codebox}[cpp]
// Chapter Filesystem/filesystem_list_files.cpp
#include <filesystem>
#include <iomanip>
#include <iostream>

namespace fs = std::filesystem;

void DisplayDirectoryTree(const fs::path& pathToScan, int level = 0) {
    for (const auto& entry : fs::directory_iterator(pathToScan)) {
        const auto filenameStr = entry.path().filename().string();
        if (entry.is_directory()) {
            std::cout << std::setw(level * 3) << "" << filenameStr << '\n';
            DisplayDirectoryTree(entry, level + 1);
        }
        else if (entry.is_regular_file()) {
            std::cout << std::setw(level * 3) << "" << filenameStr
                      << ", size " << fs::file_size(entry) << " bytes\n";
        }
        else
            std::cout << std::setw(level * 3) << "" << " [?]" << filenameStr << '\n';
    }
}

int main(int argc, char* argv[]) {
    try {
        const fs::path pathToShow{ argc >= 2 ? argv[1] : fs::current_path() };
        std::cout << "listing files in the directory: "
                  << fs::absolute(pathToShow).string() << '\n';
        DisplayDirectoryTree(pathToShow);
    }
    catch (const fs::filesystem_error& err) {
        std::cerr << "filesystem error! " << err.what() << '\n';
    }
    catch (const std::exception& ex) {
        std::cerr << "general exception: " << ex.what() << '\n';
    }
}
\end{codebox}

我们可以在临时路径 \texttt{D:\textbackslash testlist\textbackslash} 上运行这个程序，得到如下输出：

\begin{codebox}[cpp]
Running on Windows:
.\ListFiles.exe D:\testlist\
listing files in the directory: D:\testlist\
abc.txt, size 357 bytes
def.txt, size 430 bytes
ghi.txt, size 190 bytes
dirTemp
  jkl.txt, size 162 bytes
  mno.txt, size 1728 bytes
tempDir
  abc.txt, size 174 bytes
  def.txt, size 163 bytes
  tempInner
    abc.txt, size 144 bytes
    mno.txt, size 1728 bytes
    xyz.txt, size 3168 bytes
\end{codebox}

应用程序会递归地列出文件，随着每一层缩进，你可以看到我们正在进入一个新的目录。

在 Linux 上（WSL 中的 Ubuntu 18.04）运行：

\begin{codebox}[cpp]
Running on a Linux (Ubuntu 18.04 on WSL):
fenbf@FEN-NODE:/mnt/f/wsl$ ./list_files.out testList/
listing files in the directory: /mnt/f/wsl/testList/
a.txt, size 965 bytes
b.txt, size 1667 bytes
c.txt, size 1394 bytes
d.txt, size 1408 bytes
directoryTemp
  a.txt, size 1165 bytes
  b.txt, size 1601 bytes
tempDir
  a.txt, size 1502 bytes
  b.txt, size 1549 bytes
  x.txt, size 1487 bytes
\end{codebox}

现在让我们来检视这个演示的核心元素。

要使用这个库，我们必须包含相关的头文件。对于文件系统库，它是：

\begin{codebox}[cpp]
#include <filesystem>
\end{codebox}

所有的类型、函数与名字都位于 \texttt{std::filesystem} 命名空间中。为了方便，通常建立一个命名空间别名：

\begin{codebox}[cpp]
namespace fs = std::filesystem;
\end{codebox}

这样我们就可以用 \texttt{fs::path} 来引用名字，而不必写 \texttt{std::filesystem::path}。

让我们从 \texttt{main()} 函数开始，那里是应用程序逻辑的起点。

程序从命令行接收一个可选的参数。如果为空，我们就使用当前系统路径：

\begin{codebox}[cpp]
const fs::path pathToShow{ argc >= 2 ? argv[1] : fs::current_path() };
\end{codebox}

\texttt{pathToShow} 可以从字符串创建——优先使用 \texttt{argv[1]}，如果没有则 \texttt{current\_path()} 是一个返回当前系统路径的辅助自由函数。

在接下来的两行（第 28 行和第 29 行）中，我们显示迭代的起始路径。\texttt{absolute()} 函数会“展开”输入路径，并在需要时把它从相对形式转换为绝对形式。

应用程序的核心部分是 \texttt{DisplayDirectoryTree()} 函数。

在函数内部，我们使用 \texttt{directory\_iterator} 来检查目录，找出其中的文件或更多的目录。

\begin{codebox}[cpp]
for (const auto& entry : fs::directory_iterator(pathToShow))
\end{codebox}

每一次循环迭代都会产生另一个 \texttt{directory\_entry}，我们需要对它进行判断。我们可以决定是否应该递归调用该函数（当 \texttt{entry.is\_directory()} 为真时），或者如果它是一个普通文件，就只显示一些基本信息。

如你所见，我们可以访问 \texttt{path} 和 \texttt{directory entry} 的许多方法。例如，我们使用 \texttt{filename} 只返回路径中的文件名部分，从而可以显示“树”状结构。我们还调用 \texttt{fs::file\_size} 来查询文件的大小。

在简短的演示之后，现在让我们更仔细地看看 \texttt{filesystem::path}、\texttt{filesystem::directory\_entry}、辅助的非成员函数以及错误处理。

\section{路径对象}

该库的核心部分是路径对象。它包含一个路径名（pathname）——一个构成路径名称的字符串。该对象并不一定要指向文件系统中真实存在的文件。路径甚至可以是非法形式。

\myGraphic{0.9}{images/fig\_p201\_1.png}{路径的组成结构}

路径由以下元素构成：\texttt{root-name root-directory relative-path}：

\begin{itemize}
\item （可选）\texttt{root-name}：POSIX 系统没有根名。在 Windows 上，它通常是某个驱动器的名称，例如 \texttt{"C:"}
\item （可选）\texttt{root-directory}：区分相对路径与绝对路径
\item \texttt{relative-path}：
  \begin{itemize}
  \item \texttt{filename}
  \item \texttt{directory separator}
  \item \texttt{relative-path}
  \end{itemize}
\end{itemize}

我们可以用下面的示意图来说明它：

路径结构

该类实现了许多用于提取路径各部分的方法：

\begin{table}[H]
\centering
\begin{tabularx}{\textwidth}{>{\ttfamily\raggedright\arraybackslash}X >{\raggedright\arraybackslash}X}
\toprule
方法 & 说明 \\
\midrule
path::root\_name() & 返回路径的根名 \\
path::root\_directory() & 返回路径的根目录 \\
path::root\_path() & 返回路径的根路径 \\
path::relative\_path() & 返回相对于根路径的路径 \\
path::parent\_path() & 返回父路径的路径 \\
path::filename() & 返回路径中的文件名部分 \\
path::stem() & 返回路径中的主干（stem）部分 \\
path::extension() & 返回路径中的文件扩展名部分 \\
\bottomrule
\end{tabularx}
\end{table}

如果某个给定元素不存在，那么上述函数会返回一个空的路径。

还有一些方法用于查询路径的元素：

\begin{table}[H]
\centering
\begin{tabularx}{\textwidth}{>{\ttfamily\raggedright\arraybackslash}X >{\raggedright\arraybackslash}X}
\toprule
查询名 & 说明 \\
\midrule
path::has\_root\_path() & 查询路径是否拥有根 \\
path::has\_root\_name() & 查询路径是否拥有根名 \\
path::has\_root\_directory() & 检查路径是否拥有根目录 \\
path::has\_relative\_path() & 检查路径是否拥有相对路径部分 \\
path::has\_parent\_path() & 检查路径是否拥有父路径 \\
path::has\_filename() & 检查路径是否拥有文件名 \\
path::has\_stem() & 检查路径是否拥有主干部分 \\
path::has\_extension() & 检查路径是否拥有扩展名 \\
\bottomrule
\end{tabularx}
\end{table}

我们可以使用上述所有方法，并编写一个示例来显示给定路径的信息：

\begin{codebox}[cpp]
// Chapter Filesystem/filesystem_path_info.cpp
const filesystem::path testPath{...};
if (testPath.has_root_name())
    cout << "root_name() = " << testPath.root_name() << '\n';
else
    cout << "no root-name\n";
if (testPath.has_root_directory())
    cout << "root directory() = " << testPath.root_directory() << '\n';
else
    cout << "no root-directory\n";
if (testPath.has_root_path())
    cout << "root_path() = " << testPath.root_path() << '\n';
else
    cout << "no root-path\n";
if (testPath.has_relative_path())
    cout << "relative_path() = " << testPath.relative_path() << '\n';
else
    cout << "no relative-path\n";
if (testPath.has_parent_path())
    cout << "parent_path() = " << testPath.parent_path() << '\n';
else
    cout << "no parent-path\n";
if (testPath.has_filename())
    cout << "filename() = " << testPath.filename() << '\n';
else
    cout << "no filename\n";
\end{codebox}

对于文件路径（例如 \texttt{"C:\textbackslash Windows\textbackslash system.ini"}），输出如下：

\begin{codebox}[cpp]
root_name() = "C:"
root directory() = "\\"
root_path() = "C:\\"
relative_path() = "Windows\\system.ini"
parent_path() = "C:\\Windows"
filename() = "system.ini"
stem() = "system"
extension() = ".ini"
\end{codebox}

类似地，在 POSIX 系统上，对于路径 \texttt{/usr/temp/abc.txt}：

\begin{codebox}[cpp]
no root-name
root directory() = "/"
root_path() = "/"
relative_path() = "usr/temp/abc.txt"
parent_path() = "/usr/temp"
filename() = "abc.txt"
stem() = "abc"
extension() = ".txt"
\end{codebox}

还有一个小技巧可以让你遍历路径对象的各个部分。\texttt{std::filesystem::path} 为 \texttt{begin()} 和 \texttt{end()} 提供了重载，这正是你能在基于范围的 \texttt{for} 循环中使用它的原因：

\begin{codebox}[cpp]
int i = 0;
for (const auto& part : testPath)
    cout << "path part: " << i++ << " = " << part << '\n';
\end{codebox}

对于 \texttt{C:\textbackslash Windows\textbackslash system.ini} 的输出：

\begin{codebox}[cpp]
path part: 0 = C:
path part: 1 = \
path part: 2 = Windows
path part: 3 = system.ini
\end{codebox}

\subsection{路径操作}

下面是 \texttt{path} 类其他重要方法的表格：

\begin{table}[H]
\centering
\begin{tabularx}{\textwidth}{>{\ttfamily\raggedright\arraybackslash}X >{\raggedright\arraybackslash}X}
\toprule
操作 & 说明 \\
\midrule
path::append() & 用目录分隔符把一个路径追加到另一个路径之后 \\
path::concat() & 不带目录分隔符地拼接路径 \\
path::clear() & 清除所有元素，使其成为空路径 \\
path::remove\_filename() & 从路径中移除文件名部分 \\
path::replace\_filename() & 替换单个文件名部分 \\
path::replace\_extension() & 替换扩展名 \\
path::swap() & 交换两个路径 \\
path::compare() & 比较路径与另一个路径的词法表示，返回一个整数 \\
path::empty() & 检查路径是否为空 \\
\bottomrule
\end{tabularx}
\end{table}

\mySamllsectionNoContent{比较}

\texttt{path} 类有几个重载的运算符：

\texttt{==, !=, <, >, <=, =>}

以及 \texttt{path::compare()} 方法，它返回一个整数值。

所有方法都使用路径的原生格式，逐元素地进行比较。

\begin{codebox}[cpp]
fs::path p1 { "/usr/a/b/c" };
fs::path p2 { "/usr/a/b/c" };
assert(p1 == p2);
assert(p1.compare(p2) == 0);
p1 = "/usr/a/b/c";
p2 = "/usr/a/b/c/d";
assert(p1 < p2);
assert(p1.compare(p2) < 0);
\end{codebox}

在 Windows 上，我们还可以测试路径中含有根元素的情况：

\begin{codebox}[cpp]
p1 = "C:/test";
p2 = "abc/xyz"; // 没有根路径，因此它“小于”带有根的路径
assert(p1 > p2);
assert(p1.compare(p2) > 0);
\end{codebox}

或者，同样在 Windows 上，测试路径格式不同的情况：

\begin{codebox}[cpp]
fs::path p3 { "/usr/a/b/c" }; // 在 Windows 上会被转换为原生格式
fs::path p4 { "\\usr/a\\b/c" };
assert(p3 == p4);
assert(p3.compare(p4) == 0);
\end{codebox}

你可以在 Chapter Filesystem/filesystem\_compare.cpp 中把玩这些代码。

\mySamllsectionNoContent{路径拼接}

我们有两个方法可以用来拼接路径：

\begin{itemize}
\item \texttt{path::append()}——用一个目录分隔符追加路径。
\item \texttt{path::concat()}——只追加“字符串”，不带任何分隔符。
\end{itemize}

该功能也可以通过运算符 \texttt{/}、\texttt{/=}（追加）、\texttt{+} 和 \texttt{+=}（拼接）来实现。例如：

\begin{codebox}[cpp]
// append:
fs::path p1{"C:\\temp"};
p1 /= "user";
p1 /= "data";
cout << p1 << '\n';
// concat:
fs::path p2("C:\\temp\\");
p2 += "user";
p2 += "data";
cout << p2 << '\n';
\end{codebox}

输出为：

\begin{codebox}[cpp]
C:\temp\user\data
C:\temp\userdata
\end{codebox}

然而，追加路径有几条你必须留意的规则。例如，如果另一个路径是绝对路径，或者另一个路径带有根名，并且该根名与当前路径的根名不同，那么追加操作会替换掉当前路径。

\begin{codebox}[cpp]
auto resW = fs::path{"foo"} / "D:\";
// Windows
auto resP = fs::path{"foo"} / "/bar";
// POSIX
// resW 现在是 "D:\"
// resP 现在是 "/bar"
\end{codebox}

在上述情况下，\texttt{resW} 和 \texttt{resP} 都会包含第二个操作数的值，因为 \texttt{D:\textbackslash} 和 \texttt{/bar} 都包含根元素。

\mySamllsectionNoContent{流运算符}

\texttt{path} 类也实现了 \texttt{>>} 和 \texttt{<<} 运算符。

这些运算符使用 \texttt{std::quoted} 来保留正确的格式。这就是示例中的路径显示出来带有引号的原因。

在 Windows 上，这还会使运行时为原生格式的路径输出 \texttt{"\textbackslash\textbackslash"}。

例如在 POSIX 上：

\begin{codebox}[cpp]
fs::path p1 { "/usr/test/temp.xyz" };
std::cout << p1;
\end{codebox}

代码会打印 \texttt{/usr/test/temp.xyz}。

而在 Windows 上：

\begin{codebox}[cpp]
fs::path p1{ "usr/test/temp.xyz" };
fs::path p2{ "usr\\test\\temp.xyz" };
std::cout << p1 << '\n' << p2;
\end{codebox}

代码会输出：

\begin{codebox}[cpp]
"usr/test/temp.xyz"
"usr\\test\\temp.xyz"
\end{codebox}

\subsection{路径格式与转换}

文件系统库以 POSIX 为模型（例如所有基于 Unix 的系统都实现了 POSIX 标准），但也能与其他文件系统协同工作，例如 Windows。正因为如此，在以可移植的方式使用路径时有几点需要牢记。

第一点是路径格式。我们有两种核心模式：

\begin{itemize}
\item \texttt{generic}——通用格式，即标准所指定的格式（基于 POSIX 格式）
\item \texttt{native}——特定实现所使用的格式
\end{itemize}

在 POSIX 系统上，原生格式等同于通用格式。但在 Windows 上则不同。

格式的主要区别在于，Windows 使用反斜杠（\texttt{\textbackslash}）而不是斜杠（\texttt{/}）。另一点是 Windows 拥有根目录——例如 \texttt{C:}、\texttt{D:} 或其他盘符。

还有一点重要的是用于保存路径元素的字符串类型。在 POSIX 上是 \texttt{char}（以及 \texttt{std::string}），而在 Windows 上则是 \texttt{wchar\_t} 和 \texttt{std::wstring}。\texttt{path} 类型定义了 \texttt{path::value\_type} 和 \texttt{string\_type}（定义为 \texttt{std::basic\_string<value\_type>}）来暴露这些属性。

\texttt{path} 类有几个方法，让你可以选用最匹配的格式。

如果你想使用原生格式，可以使用：

\begin{table}[H]
\centering
\begin{tabularx}{\textwidth}{>{\ttfamily\raggedright\arraybackslash}X >{\raggedright\arraybackslash}X}
\toprule
操作 & 说明 \\
\midrule
path::c\_str() & 返回 \texttt{value\_type*} \\
path::native() & 返回 \texttt{string\_type\&} \\
\bottomrule
\end{tabularx}
\end{table}

还有许多方法会把原生格式进行转换：

\begin{table}[H]
\centering
\begin{tabularx}{\textwidth}{>{\ttfamily\raggedright\arraybackslash}X >{\raggedright\arraybackslash}X}
\toprule
操作 & 说明 \\
\midrule
path::string() & 转换为 \texttt{string} \\
path::wstring() & 转换为 \texttt{wstring} \\
path::u8string() & 转换为 \texttt{u8string} \\
path::u16string() & 转换为 \texttt{u16string} \\
path::u32string() & 转换为 \texttt{u32string} \\
\bottomrule
\end{tabularx}
\end{table}

由于 Windows 使用 \texttt{wchar\_t} 作为路径的底层类型，因此你需要留意向 \texttt{char} 的“隐藏”转换。

\section{目录项与目录遍历}

虽然 \texttt{path} 类代表存在或不存在的文件或路径，但我们还有另一个更具体的对象：它就是 \texttt{directory\_entry} 对象。该对象指向真实存在的文件和目录，并且通常通过文件系统迭代器获得。

此外，实现被鼓励去缓存额外的文件属性。这样一来就能减少系统调用的次数。

\subsection{使用目录迭代器遍历路径}

你可以使用两个可用的迭代器来遍历路径：

\begin{itemize}
\item \texttt{directory\_iterator}——在单个目录内迭代，输入迭代器。
\item \texttt{recursive\_directory\_iterator}——递归迭代，输入迭代器。
\end{itemize}

在两种方式中，被访问文件名的顺序都是未指定的，每个目录项只会被访问一次。

如果在目录迭代器创建之后，有文件或目录被删除或添加到目录树中，那么这种变化是否会透过迭代器被观察到是未指定的。

在两种迭代器中，目录 \texttt{.} 和 \texttt{..} 都会被跳过。

你可以使用下面的模式来遍历一个目录：

\begin{codebox}[cpp]
for (auto const & entry : fs::directory_iterator(pathToShow))
{
    ...
}
\end{codebox}

或者使用另一种方式，借助算法，你还可以在此基础上过滤路径：

\begin{codebox}[cpp]
std::filesystem::recursive_directory_iterator dirpos{ inPath };
std::copy_if(begin(dirpos), end(dirpos), std::back_inserter(paths),
    some_predicate);
\end{codebox}

\texttt{some\_predicate} 是一个谓词，根据某个给定的 \texttt{directory\_entry} 对象是否匹配我们的过滤器，返回 \texttt{true} 或 \texttt{false}。

\subsection{directory\_entry 的方法}

下面是 \texttt{directory\_entry} 方法的列表：

\begin{table}[H]
\centering
\begin{tabularx}{\textwidth}{>{\ttfamily\raggedright\arraybackslash}X >{\raggedright\arraybackslash}X}
\toprule
操作 & 说明 \\
\midrule
directory\_entry::assign() & 替换项内的路径，并调用 \texttt{refresh()} 以更新缓存的属性 \\
directory\_entry::replace\_filename() & 替换项内的文件名，并调用 \texttt{refresh()} 以更新缓存的属性 \\
directory\_entry::refresh() & 更新文件的缓存属性 \\
directory\_entry::exists() & 检查目录项是否指向真实存在的文件系统对象 \\
directory\_entry::is\_block\_file() & 如果文件项是块设备文件则返回 \texttt{true} \\
directory\_entry::is\_character\_file() & 如果文件项是字符设备文件则返回 \texttt{true} \\
directory\_entry::is\_directory() & 如果文件项是目录则返回 \texttt{true} \\
directory\_entry::is\_fifo() & 如果文件项指向一个命名管道则返回 \texttt{true} \\
directory\_entry::is\_other() & 如果文件项指向其他文件类型则返回 \texttt{true} \\
directory\_entry::is\_regular\_file() & 如果文件项是普通文件则返回 \texttt{true} \\
directory\_entry::is\_socket() & 如果文件项是一个命名 IPC 套接字则返回 \texttt{true} \\
directory\_entry::is\_symlink() & 如果文件项是一个符号链接则返回 \texttt{true} \\
directory\_entry::file\_size() & 返回该目录项所指向文件的大小 \\
directory\_entry::hard\_link\_count() & 返回指向该文件的硬链接数量 \\
directory\_entry::last\_write\_time() & 获取或设置文件的最后写入时间 \\
directory\_entry::status() & 返回该目录项所指定文件的状态 \\
directory\_entry::symlink\_status() & 返回该目录项所指定文件的符号链接状态 \\
\bottomrule
\end{tabularx}
\end{table}

\section{辅助函数}

到目前为止，我们已经介绍了文件系统的三个元素：\texttt{path} 类、\texttt{directory\_entry} 和目录迭代器。该库还提供了一组非成员函数。

\mySamllsectionNoContent{查询函数}

\begin{table}[H]
\centering
\begin{tabularx}{\textwidth}{>{\ttfamily\raggedright\arraybackslash}X >{\raggedright\arraybackslash}X}
\toprule
函数 & 说明 \\
\midrule
filesystem::is\_block\_file() & 检查给定路径是否指向块设备 \\
filesystem::is\_character\_file() & 检查给定路径是否指向字符设备 \\
filesystem::is\_directory() & 检查给定路径是否指向目录 \\
filesystem::is\_empty() & 检查给定路径是否指向空文件或空目录 \\
filesystem::is\_fifo() & 检查给定路径是否指向命名管道 \\
filesystem::is\_other() & 检查参数是否指向另一个文件 \\
filesystem::is\_regular\_file() & 检查参数是否指向普通文件 \\
filesystem::is\_socket() & 检查参数是否指向命名 IPC 套接字 \\
filesystem::is\_symlink() & 检查参数是否指向符号链接 \\
filesystem::status\_known() & 检查文件状态是否已知 \\
filesystem::exists() & 检查路径是否指向真实存在的文件系统对象 \\
filesystem::file\_size() & 返回文件的大小 \\
filesystem::last\_write\_time() & 获取或设置最后一次数据修改的时间 \\
\bottomrule
\end{tabularx}
\end{table}

\mySamllsectionNoContent{与路径相关}

\begin{table}[H]
\centering
\begin{tabularx}{\textwidth}{>{\ttfamily\raggedright\arraybackslash}X >{\raggedright\arraybackslash}X}
\toprule
函数名 & 说明 \\
\midrule
filesystem::absolute() & 组合出一个绝对路径 \\
filesystem::canonical() & 组合出一个规范路径 \\
filesystem::weakly\_canonical() & 组合出一个弱规范路径 \\
filesystem::relative() & 组合出一个相对路径 \\
filesystem::proximate() & 组合出一个邻近路径 \\
filesystem::current\_path() & 返回或设置当前工作目录 \\
filesystem::equivalent() & 检查两个路径是否指向同一个文件系统对象 \\
\bottomrule
\end{tabularx}
\end{table}

\mySamllsectionNoContent{目录与文件管理}

\begin{table}[H]
\centering
\begin{tabularx}{\textwidth}{>{\ttfamily\raggedright\arraybackslash}X >{\raggedright\arraybackslash}X}
\toprule
函数名 & 说明 \\
\midrule
filesystem::copy() & 复制文件或目录 \\
filesystem::copy\_file() & 复制文件内容 \\
filesystem::copy\_symlink() & 复制一个符号链接 \\
filesystem::create\_directory(), filesystem::create\_directories() & 创建新目录 \\
filesystem::create\_hard\_link() & 创建硬链接 \\
filesystem::create\_symlink(), filesystem::create\_directory\_symlink() & 创建符号链接 \\
filesystem::hard\_link\_count() & 返回指向特定文件的硬链接数量 \\
filesystem::permissions() & 修改文件访问权限 \\
filesystem::read\_symlink() & 获取符号链接的目标 \\
filesystem::remove(), filesystem::remove\_all() & 删除单个文件，或递归删除整个目录及其全部内容 \\
filesystem::rename() & 移动或重命名文件或目录 \\
filesystem::resize\_file() & 通过截断或填零来改变普通文件的大小 \\
filesystem::space() & 确定文件系统上可用的空闲空间 \\
filesystem::status(), filesystem::symlink\_status() & 确定文件属性、确定文件属性、检查符号链接目标 \\
filesystem::temp\_directory\_path() & 返回一个适合作为临时文件的目录 \\
\bottomrule
\end{tabularx}
\end{table}

\subsection{获取与显示文件时间}

在 C++17 中，关于 \texttt{last\_write\_time()} 的值有一点不太方便。

我们有一个自由函数，以及 \texttt{directory\_entry} 中的一个方法。它们都返回 \texttt{file\_time\_type}，目前它被定义为：

\begin{codebox}[cpp]
using file_time_type = std::chrono::time_point</*trivial-clock*/>;
\end{codebox}

来自标准 30.10.25 头文件 \texttt{<filesystem>} 的概要：

\texttt{trivial-clock} 是一个由实现定义的类型，它满足 TrivialClock 要求，并且有能力表示并测量文件时间值。实现应当确保文件时间类型的分辨率与范围，能够反映操作系统相关的文件时间值的分辨率与范围。

换句话说，这取决于具体的实现。

例如在 GCC/Clang 的 STL 中，文件时间是基于 \texttt{chrono::system\_clock} 实现的，而在 MSVC 中它则是一个平台相关的时钟。

关于 Visual Studio 中的实现决策，这里有更多细节：\texttt{std::filesystem::file\_time\_type} 不允许轻松转换为 \texttt{time\_t}\footnote{\url{https://developercommunity.visualstudio.com/content/problem/251213/stdfilesystemfile-time-type-does-not-allow-easy-co.html}}。

这种情况可能很快就会改善，因为在 C++20 中我们将获得 \texttt{std::chrono::file\_clock} 以及时钟之间的转换例程。参见 P0355\footnote{\url{https://wg21.link/p0355}}（已加入 C++20）。

让我们看一些代码。

\begin{codebox}[cpp]
auto filetime = fs::last_write_time(myPath);
const auto toNow = fs::file_time_type::clock::now() - filetime;
const auto elapsedSec = duration_cast<seconds>(toNow).count();
// skipped std::chrono prefix for duration_cast and seconds
\end{codebox}

上面的代码给出了一种计算自上次更新以来所经过秒数的方法。不过，这不如显示一个真实日期来得有用。

在 POSIX 上（GCC 和 Clang 的实现）你可以轻松地把文件时间转换为 \texttt{system\_clock}，然后再获取 \texttt{std::time\_t}：

\begin{codebox}[cpp]
auto filetime = fs::last_write_time(myPath);
std::time_t convfiletime = std::chrono::system_clock::to_time_t(filetime);
std::cout << "Updated: " << std::ctime(&convfiletime) << '\n';
\end{codebox}

在 MSVC 中这段代码无法编译。不过有一个保证：\texttt{file\_time\_type} 可以与接收 \texttt{FILETIME} 的原生操作系统函数一起使用。因此我们可以编写如下代码来解决这个问题：

\begin{codebox}[cpp]
auto filetime = fs::last_write_time(myPath);
FILETIME ft;
memcpy(&ft, &filetime, sizeof(FILETIME));
SYSTEMTIME
stSystemTime;
if (FileTimeToSystemTime(&ft, &stSystemTime)) {
    // use stSystemTime.wYear, stSystemTime.wMonth, stSystemTime.wDay, ...
}
\end{codebox}

请参阅 Chapter Filesystem/filesystem\_list\_files\_info.cpp 查看示例。

\subsection{文件权限}

在上面的表格中你可能已经注意到与文件权限相关的函数。我们有两个主要的函数：

\begin{itemize}
\item \texttt{std::filesystem::status()} 和
\item \texttt{std::filesystem::permissions()}
\end{itemize}

第一个函数返回 \texttt{file\_status}，其中包含关于文件类型以及其权限的信息。

你可以使用第二个函数来修改文件权限。例如，把一个文件改成只读。

文件权限——\texttt{std::filesystem::perms}——它是一个枚举类，表示如下的取值：

\begin{table}[H]
\centering
\begin{tabularx}{\textwidth}{>{\ttfamily\raggedright\arraybackslash}X c >{\ttfamily}c >{\raggedright\arraybackslash}X}
\toprule
名称 & 值（八进制） & POSIX 宏 & 说明 \\
\midrule
none & & & 没有为文件设置任何权限 \\
owner\_read & 0400 & S\_IRUSR & 所有者读权限 \\
owner\_write & 0200 & S\_IWUSR & 所有者写权限 \\
owner\_exec & 0100 & S\_IXUSR & 所有者执行/搜索权限 \\
owner\_all & 0700 & S\_IRWXU & 所有者的读、写、执行/搜索权限 \\
group\_read & 0040 & S\_IRGRP & 组读权限 \\
group\_write & 0020 & S\_IWGRP & 组写权限 \\
group\_exec & 0010 & S\_IXGRP & 组执行/搜索权限 \\
group\_all & 0070 & S\_IRWXG & 组的读、写、执行/搜索权限 \\
others\_read & 0004 & S\_IROTH & 其他用户读权限 \\
others\_write & 0002 & S\_IWOTH & 其他用户写权限 \\
others\_exec & 0001 & S\_IXOTH & 其他用户执行/搜索权限 \\
others\_all & 0007 & S\_IRWXO & 其他用户的读、写、执行/搜索权限 \\
all & 0777 & owner\_all | group\_all | others\_all & 所有用户的读、写、执行/搜索权限 \\
set\_uid & 04000 & S\_ISUID & 执行时设置用户 ID \\
set\_gid & 02000 & S\_ISGID & 执行时设置组 ID \\
sticky\_bit & 01000 & S\_ISVTX & 取决于操作系统 \\
mask & 07777 & all | set\_uid | set\_gid | sticky\_bit & \\
unknown & 0xFFFF & & 权限未知 \\
\bottomrule
\end{tabularx}
\end{table}

下面是一个简短的代码示例，演示如何打印文件权限：

\begin{codebox}[cpp]
// Chapter Filesystem/filesystem_permissions.cpp
std::ostream& operator<< (std::ostream& stream, fs::perms p)
{
    stream << "owner: "
           << ((p & fs::perms::owner_read)
               != fs::perms::none ? "r" : "-")
           << ((p & fs::perms::owner_write) != fs::perms::none ? "w" : "-")
           << ((p & fs::perms::owner_exec)
               != fs::perms::none ? "x" : "-");
    stream << " group: "
           << ((p & fs::perms::group_read)
               != fs::perms::none ? "r" : "-")
           << ((p & fs::perms::group_write) != fs::perms::none ? "w" : "-")
           << ((p & fs::perms::group_exec)
               != fs::perms::none ? "x" : "-");
    stream << " others: "
           << ((p & fs::perms::others_read)
               != fs::perms::none ? "r" : "-")
           << ((p & fs::perms::others_write) != fs::perms::none ? "w" : "-")
           << ((p & fs::perms::others_exec)
               != fs::perms::none ? "x" : "-");
    return stream;
}
\end{codebox}

你可以像下面这样使用上面的 \texttt{operator<<} 实现：

\begin{codebox}[cpp]
std::cout << "perms: " << fs::status("myFile.txt").permissions() << '\n';
\end{codebox}

\mySamllsectionNoContent{设置权限}

要修改权限，可以使用如下代码：

\begin{codebox}[cpp]
std::cout << "after creation: " << fs::status(sTempName).permissions() << '\n';
fs::permissions(sTempName, fs::perms::owner_read, fs::perm_options::remove);
std::cout << "after change: " << fs::status(sTempName).permissions() << '\n';
\end{codebox}

\texttt{std::filesystem::permissions} 是一个函数，它接收一个路径，然后一个标志以及“动作”参数。

\texttt{fs::perm\_options} 有三种模式：

\begin{itemize}
\item \texttt{replace}——你传入的权限标志会替换现有状态。这是该参数的默认值。
\item \texttt{add}——权限标志会与现有状态做按位或（OR）操作。
\item \texttt{remove}——权限会被替换为“取反后的参数与当前权限的按位与（AND）”。
\item \texttt{nofollow}——权限会在符号链接本身上修改，而不是它所解析到的文件上。
\end{itemize}

例如：

\begin{codebox}[cpp]
// remove "owner_read"
fs::permissions(myPath, fs::perms::owner_read, fs::perm_options::remove);
// add "owner_read"
fs::permissions(myPath, fs::perms::owner_read, fs::perm_options::add);
// replace and set "owner_all":
fs::permissions(myPath, fs::perms::owner_all); // replace is default param
\end{codebox}

\mySamllsectionNoContent{Windows 注意事项}

Windows 不是一个 POSIX 系统，它不会把 POSIX 文件权限映射到自己的体系上。对于 \texttt{std::filesystem}，它只支持两种模式：只读和全部。

来自微软文档（Microsoft Docs）的文件系统文档\footnote{\url{https://docs.microsoft.com/en-us/cpp/standard-library/filesystem-enumerations?view=vs-2017}}：

受支持的取值本质上是“只读”（readonly）和全部（all）。对于一个只读文件，所有的 \texttt{*\_write} 位都不会被设置。否则就会设置全部位（0777）。

因此，遗憾的是，如果你想在 Windows 上修改文件权限，你的选择非常有限。

\section{错误处理与文件竞态}

到目前为止，本章的示例都使用异常处理来报告错误。文件系统 API 还配备了会输出错误码的函数与方法重载。你可以自行决定是要异常还是要错误码。

例如我们为 \texttt{file\_size} 提供了两个重载：

\begin{codebox}[cpp]
uintmax_t file_size(const path& p);
uintmax_t file_size(const path& p, error_code& ec) noexcept;
\end{codebox}

第二个重载可以这样使用：

\begin{codebox}[cpp]
const std::filesystem::path testPath("C:\test.txt");
std::error_code ec{};
auto size = std::filesystem::file_size(testPath, ec);
if (ec == std::error_code{})
    std::cout << "size: " << size << '\n';
else
    std::cout << "error when accessing test file, size is: "
              << size << " message: " << ec.message() << '\n';
\end{codebox}

\texttt{file\_size} 接收一个额外的输出参数——\texttt{error\_code}，并在发生状况时设置一个值。如果操作成功，\texttt{ec} 会被值初始化。

\subsection{文件竞态}

需要重点指出的是，当发生文件竞态时，可能会出现未定义行为。

来自 30.10.2.3 文件系统竞态行为：

\begin{myQuote}
如果对本库函数的调用引入了文件系统竞态，则行为未定义。
\end{myQuote}

以及来自 30.10.9 文件系统竞态：

\begin{myQuote}
当多个线程、进程或计算机在文件系统中交错地访问并修改同一个对象时，所出现的情况。
\end{myQuote}

\section{示例}

在本节中，我们将分析几个使用 \texttt{std::filesystem} 的示例。我们会从一个简单的情形入手——把文件加载到字符串中，然后探究如何创建目录，再使用 \texttt{std::regex} 来过滤文件名。

演示样例位于 Chapter Filesystem/filesystem\_list\_files.cpp，而它的扩展版本（会显示文件时间与大小）位于 Chapter Filesystem/filesystem\_list\_files\_info.cpp。

要了解更多用例，你还可以阅读《如何并行化 CSV 读取器》一章，其中 \texttt{std::filesystem} 是查找 CSV 文件的关键元素。

\subsection{将文件加载到字符串中}

第一个示例展示了一个很有说服力的情形：我们利用 \texttt{std::filesystem} 与大小相关的函数，为文件内容构建一个缓冲区。

代码如下：

\begin{codebox}[cpp]
// Chapter Filesystem/filesystem_load_string.cpp
[[nodiscard]] std::string GetFileContents(const fs::path& filePath)
{
    std::ifstream inFile{ filePath, std::ios::in | std::ios::binary };
    if (!inFile)
        throw std::runtime_error("Cannot open " + filePath.string());
    const auto fsize = fs::file_size(filePath);
    if (fsize > std::numeric_limits<size_t>::max())
        throw std::runtime_error("file is too large to fit into size_t! "
                                 + filePath.string());
    std::string str(static_cast<size_t>(fsize), 0);
    inFile.read(str.data(), str.size());
    if (!inFile)
        throw std::runtime_error("Could not read the full contents from "
                                 + filePath.string());
    return str;
}
\end{codebox}

在 C++17 之前，要获取文件大小，你通常会把文件指针重新定位到末尾，然后再读取一次位置。例如：

\begin{codebox}[cpp]
ifstream testFile("test.file", ios::binary);
const
auto begin = myfile.tellg();
testFile.seekg (0, ios::end);
const auto end = testFile.tellg();
const auto fsize =
(end-begin);
\end{codebox}

你也可以带上 \texttt{ios::ate} 标志打开文件，这样文件指针会自动定位到末尾。

然而，上述所有方法都需要打开一个文件，而使用 \texttt{std::filesystem} 代码则要短得多，并且我们只需要读取文件属性。

更重要的是，\texttt{std::filesystem} 技术所需的访问权限“更低”，因为你只需要父目录的读权限，而不需要“文件读”权限。

如果你使用 \texttt{std::filesystem::directory\_entry} 的方法，那么文件大小有可能来自缓存。

\subsection{创建目录}

在第二个示例中，我们将构建 N 个目录，每个目录中包含 M 个文件。

\texttt{main()} 的核心部分：

\begin{codebox}[cpp]
// Chapter Filesystem/filesystem_build_temp.cpp
const fs::path startingPath{ argc >= 2 ? argv[1] : fs::current_path() };
const std::string strTempName{ argc >= 3 ? argv[2] : "temp" };
const int numDirectories{ argc >= 4 ? std::stoi(argv[3]) : 4 };
const int numFiles{ argc >= 5 ? std::stoi(argv[4]) : 4 };
if (numDirectories < 0 || numFiles < 0)
    throw std::runtime_error("negative input numbers...");
const fs::path tempPath = startingPath / strTempName;
CreateTempData(tempPath, numDirectories, numFiles);
\end{codebox}

以及 \texttt{CreateTempData()} 函数：

\begin{codebox}[cpp]
// Chapter Filesystem/filesystem_build_temp.cpp
std::vector<fs::path> GeneratePathNames(const fs::path& tempPath,
                                        unsigned num)
{
    std::vector<fs::path> outPaths{ num, tempPath };
    for (auto& dirName : outPaths)
    {
        // use pointer value to generate unique name...
        const auto addr = reinterpret_cast<uintptr_t>(&dirName);
        dirName /= std::string("tt") + std::to_string(addr);
    }
    return outPaths;
}

void CreateTempFiles(const fs::path& dir, unsigned numFiles)
{
    auto files = GeneratePathNames(dir, numFiles);
    for (auto &oneFile : files)
    {
        std::ofstream entry(oneFile.replace_extension(".txt"));
        entry << "Hello World";
    }
}

void CreateTempData(const fs::path& tempPath, unsigned numDirectories,
                    unsigned numFiles)
{
    fs::create_directory(tempPath);
    auto dirPaths = GeneratePathNames(tempPath, numDirectories);
    for (auto& dir : dirPaths)
    {
        if (fs::create_directory(dir))
            CreateTempFiles(dir, numFiles);
    }
}
\end{codebox}

在 \texttt{CreateTempData()} 中，我们首先创建文件夹结构的根。然后在 \texttt{GeneratePathNames()} 中生成路径名。每个路径名都由指针地址构成。这种方式应该能给我们提供足够唯一的名字。当我们拥有了一个唯一路径的向量之后，我们再生成另一个用于文件的路径向量。在这种情况下，我们使用 \texttt{.txt} 扩展名。使用 \texttt{fs::create\_directory} 创建目录，而创建文件则可以使用标准的流对象。

如果我们用如下参数运行应用程序：“. temp 2 4”，它会创建如下的目录结构：

\begin{codebox}[cpp]
temp
  tt22325368
    tt22283456.txt, size 11 bytes
    tt22283484.txt, size 11 bytes
    tt22283512.txt, size 11 bytes
    tt22283540.txt, size 11 bytes
  tt22325396
    tt22283456.txt, size 11 bytes
    tt22283484.txt, size 11 bytes
    tt22283512.txt, size 11 bytes
    tt22283540.txt, size 11 bytes
\end{codebox}

\subsection{使用正则表达式过滤文件}

本章的最后一个示例将借助 \texttt{std::regex}（自 C++11 起可用）来过滤文件名。

\texttt{main()} 的核心部分：

\begin{codebox}[cpp]
// Chapter Filesystem/filesystem_filter_files.cpp
const fs::path pathToShow{ argc >= 2 ? argv[1] : fs::current_path() };
const std::regex reg(argc >= 3 ? argv[2] : "");
auto files = CollectFiles(pathToShow);
std::sort(files.begin(), files.end());
for (auto& entry : files)
{
    const auto strFileName = entry.mPath.filename().string();
    if (std::regex_match(strFileName, reg))
        std::cout << strFileName << "\tsize: " << entry.mSize << '\n';
}
\end{codebox}

应用程序会收集给定目录（递归地）中的所有文件。随后，如果某个文件项与正则匹配，就显示出来。

\texttt{CollectFiles()} 函数使用递归目录迭代器来查找所有文件，并为每个文件构建基本信息。在 \texttt{main()} 中我们按大小对每个文件排序。

\begin{codebox}[cpp]
// Chapter Filesystem/filesystem_filter_files.cpp
struct FileEntry
{
    fs::path mPath;
    uintmax_t mSize{ 0 };
    static FileEntry Create(const fs::path& filePath) {
        return FileEntry{ filePath, fs::file_size(filePath) };
    }
    friend bool operator < (const FileEntry& a, const FileEntry& b) noexcept {
        return a.mSize < b.mSize;
    }
};

std::vector<FileEntry> CollectFiles(const fs::path& inPath)
{
    std::vector<fs::path> paths;
    if (fs::exists(inPath) && fs::is_directory(inPath))
    {
        std::filesystem::recursive_directory_iterator dirpos{ inPath };
        std::copy_if(begin(dirpos), end(dirpos), std::back_inserter(paths),
            [](const fs::directory_entry& entry) {
                return entry.is_regular_file();
            }
        );
    }
    std::vector<FileEntry> files(paths.size());
    std::transform(paths.cbegin(), paths.cend(), files.begin(), FileEntry::Create);
    return files;
}
\end{codebox}

在 \texttt{CollectFiles} 中，我们使用一个递归迭代器，然后借助 \texttt{std::copy\_if} 只过滤出普通文件。之后，一旦文件被收集完毕，我们就创建 \texttt{FileEntry} 对象的输出向量。\texttt{FileEntry::Create()} 会初始化对象，并同时获取文件的大小。

例如，如果我们用如下参数运行应用程序“temp .*.txt”，我们就会查找目录中所有的 txt 文件。

\begin{codebox}[cpp]
.\FilterFiles.exe temp .*.txt
tt22283456.txt
  size: 11
tt22283484.txt
  size: 11
tt22283512.txt
  size: 11
tt22283540.txt
  size: 11
tt22283456.txt
  size: 11
tt22283484.txt
  size: 11
tt22283512.txt
  size: 11
tt22283540.txt
  size: 11
\end{codebox}

\mySamllsectionNoContent{优化与代码清理}

\texttt{CollectFiles()} 函数会遍历一个目录，然后把普通文件的路径输出到一个向量中。之后该函数创建 \texttt{FileEntry} 对象，并将其返回到一个独立的向量中。

问题在于，我们没有充分利用扫描期间所拥有的 \texttt{filesystem::directory\_entry} 对象的全部可能。例如，\texttt{filesystem::directory\_entry::file\_size} 方法会比 \texttt{filesystem::file\_size} 快得多，因为 \texttt{directory\_entry} 通常会把文件属性保留在缓存中。

另一个可以优化的点是那个临时的路径向量。我们可以通过基于范围的 \texttt{for} 循环来跳过它。

下面是最终的代码：

\begin{codebox}[cpp]
// Chapter Filesystem/filesystem_filter_files.cpp
std::vector<FileEntry> CollectFilesOpt(const fs::path& inPath)
{
    std::vector<FileEntry> files;
    if (fs::exists(inPath) && fs::is_directory(inPath))
    {
        for (const auto& entry : fs::recursive_directory_iterator{ inPath })
        {
            if (entry.is_regular_file())
                files.push_back({ entry, entry.file_size() });
        }
    }
    return files;
}
\end{codebox}

\section{本章小结}

在本章中，我们深入探究了 C++17 最大的新增特性之一：\texttt{std::filesystem}。你看到了该库的核心元素：\texttt{path} 类、\texttt{directory\_entry} 和迭代器，以及大量辅助自由函数。

贯穿全章，我们还探究了许多示例：从简单的情形（如拼接路径、获取文件大小、遍历目录）到更加复杂的情形（如使用正则过滤，或创建临时目录结构）。

你现在应该已经具备了关于 \texttt{std::filesystem} 的扎实知识，并准备好自己去探索这个库了。

\texttt{std::filesystem} 的完整实现在论文 P0218：为 C++17 采纳文件系统 TS\footnote{\url{https://wg21.link/p0218}} 中描述。还有其他一些更新，例如 P0317：目录项缓存\footnote{\url{https://wg21.link/p0317}}，P0430——在非 POSIX 类操作系统上的文件系统库\footnote{\url{https://wg21.link/p0430}}，P0492R2——C++17 国家机构意见的解决方案\footnote{\url{https://wg21.link/p0492}}，P0392——通过文件系统路径适配 string\_view\footnote{\url{https://wg21.link/p0392r0}}。你可以在 C++17 草案 N4687\footnote{\url{https://wg21.link/n4687}} 的“filesystem”一节（30.10）中找到最终规范，或者在这个在线地址 timsong-cpp/filesystems\footnote{\url{https://timsong-cpp.github.io/cppwp/n4659/filesystems}} 下找到。

\section{编译器支持}

\paragraph{GCC/libstdc++}

该库在 8.0 版本中加入——参见这个提交——实现 C++17 文件系统库\footnote{\url{https://github.com/gcc-mirror/gcc/commit/3b90ed62fb848046ed7ddef07df7c806e7f3fadb}}。不过，从 GCC 5.3 起你就可以把玩实验性版本——即 TS 实现。

在这两种情况（完整支持或实验性）下，你都必须链接 \texttt{-lstdc++fs}，因为该库位于主标准库模块之外。

例如要编译示例代码——Chapter Filesystem/filesystem\_list\_files.cpp——你应该编写如下命令：

\begin{codebox}[text]
g++ -std=c++17 -O2 -Wall -Werror -pedantic filesystem_list_files.cpp -lstdc++fs
\end{codebox}

\paragraph{Clang/libc++}

对 \texttt{<filesystem>} 的支持最近在 7.0 版本中实现，你可以看到这个提交——实现 filesystem\footnote{\url{https://github.com/llvm-mirror/libcxx/commit/a0866c5fb5020d15c69deda94d32a7f982b88dc9}}。

从 Clang 3.9 起你就可以开始把玩实验性版本，即 TS 实现。

与 GCC 类似，你必须链接 \texttt{libc++fs.a}。

\paragraph{Visual Studio}

\texttt{<filesystem>} 的完整实现在 Visual Studio 2017 15.7 中加入。

在 15.7 之前，你可以在更早的版本中使用 \texttt{<experimental/filesystem>}。实验性实现甚至在 Visual Studio 2012 中就已可用，并在之后的每个版本中逐步改进。

\subsection{编译器支持汇总}

\begin{table}[H]
\centering
\begin{tabularx}{\textwidth}{X c c c}
\toprule
特性 & GCC & Clang & MSVC \\
\midrule
Filesystem & 8.0 & 7.0 & VS 2017 15.7 \\
\bottomrule
\end{tabularx}
\end{table}
